\SourceLeaf{167}

Eher könnte als zulässig erscheinen die Einrichtung, dass für die Witwen zweiter Klasse, neben gleichmässiger Theilnahme an der Zusatzdividende, der fixe Theil der Pension niedriger als 200 Rthl. z. B. zu 150 Rthl. festgesetzt würde. Erwägt man jedoch,

dass die Einbusse an Beiträgen sogleich anfängt, sobald die zweite Klasse constituirt ist, und dann fortwährend zunimmt, je mehr neue ausserordentliche Professoren ernannt werden;

dass jetzt unter den 51 Mitgliedern der Witwenkasse 20 ausserordentliche Professoren sich befinden, und folglich, wenn auch in Zukunft ein ähnliches Verhältniss bleibt, von der im Hauptplane vorausgesetzten Beitragserhöhung ein nach und nach bis zu \(\tfrac{2}{5}\) anwachsender Theil ausbleibt; indess dagegen

der Ersatz aus eintretenden geringern Pensionirungen aller Wahrscheinlichkeit nach erst nach langer Zeit anfangen, und bei der sehr geringen Mortalität, die demjenigen Alter zukommt, in dem die Mehrzahl der ausserordentlichen Professoren zu stehen pflegt, auch, durchschnittlich, selten sein wird

so bleibt es sehr problematisch, ob die Bestimmung des festen Theils der Pension zu 150 Rthl. für die zweite Klasse zureichen würde, der Kasse auch nur den vollen Ersatz für den Verlust durch die geringern Beiträge zu gewähren. Noch viel weniger aber dürfte man eine solche Maassregel als entschieden zum \textit{Vortheil} der Kasse gereichend annehmen. [Dazu wäre man nur berechtigt, wenn man entweder die Beiträge in der zweiten Klasse eben so hoch wie in der ersten beibehielte, oder die Pension noch erheblich unter 150 Rthl. herabsetzte; allein das eine wie das andre würde so sehr wie eine Härte erscheinen, dass die Commission sich nicht dafür erklären kann.]

Eine andere Maassregel, durch welche zuweilen der Kasse einiger Vortheil zuwachsen könnte, wäre die Aufhebung der im 10. Artikel des Regulativs enthaltenen Bestimmung, nach welcher die Witwenpension durch eine Wiederverheirathung der Witwe ganz erlischt. Der Zweck dieser Verordnung kann nur gewesen sein, dass man in der Voraussetzung, solche Fälle würden öfters vorkommen, der Kasse einen Gewinn hat zuwenden wollen. Allein dieser Zweck wird so gut wie ganz verfehlt, da der Erfahrung zufolge Wiederverheirathungen der Witwen unter diesen Umständen etwas fast Unerhörtes sind. (Es ist schon in der Denkschrift bemerkt, dass in mehr als 100 Jahren nur Ein solcher Fall vorgekommen ist.) Zweckmässiger ist ohne Zweifel die bei andern Witwenkassen bestehende Einrichtung, dass die Witwenpension während eines zweiten Ehestandes nur ruhet, aber wieder auflebt, wenn die Wiederverheirathete zum zweiten Male Witwe wird. Man vergleiche die Statuten der Hof- und Civil-Diener Witwenkasse §. 25, der Prediger-Witwenkasse §. 24, der Schullehrer-Witwenkasse §. 22. Der Fall, wo der zweite Ehemann wieder ein hiesiger Professor ist, würde wohl ausgenommen werden müssen, da es unzulässig scheint, aus unsrer Witwenkasse Einer Witwe (oder Einer Familie) zwei Pensionen zu gewähren. Dagegen aber dürfte es rathsam sein, von weitern Beschränkungen der Reviviscenz Umgang zu nehmen, und namentlich die volle Professorenwitwenpension auch für den Fall wieder zuzuführen, wo die Witwe durch ihre zweite Verheirathung z. B. mit einem Prediger oder andern Staatsdiener eine Pension aus einer \textit{andern} öffentlichen Kasse erhielte. Man muss nemlich erwägen, dass von einer Abänderung des bisherigen Status nur in so weit eine Wirkung erwartet werden kann, als die Abänderung nicht wieder durch Ausnahme-Verfügungen aufgehoben ist. Dass übrigens in dem Falle, wo aus der ersten Ehe Kinder unter 20 Jahren vorhanden sind, die Waisenpension auch während der zweiten Ehe in demselben Maasse wie bisher fortdauern musste, versteht sich von selbst.

\Needspace{5\baselineskip}
\pdfbookmark[2]{Note zum Commissionsberichte.}{witwen-iii-note}
\begin{center}\textit{Note zum Commissionsberichte.}\end{center}
\noindent\hspace*{\parindent}Wenn die Ausdrücke Bilanz und Deficit in Beziehung auf das Finanzbudget eines Staats gebraucht werden, so versteht man unter ersterer die Vergleichung der Ausgaben und Einnahmen, wie sie für Ein Jahr, oder für eine kleine Anzahl von Jahren, die eine Finanzperiode bilden, nach präsumtiver Veran%
